
\begin{proposition}[A fixed punctured torus family has infinitely many zero-Euler fillings with essential projective elliptic components of coherent Euler characteristic one]\label{cdp:st:C.2.UN.27}
For every integer $r\ge2$, the compact construction
of \ref{cdp:st:C.2.UN.25} has a version whose distinguished smooth
component $M_r$ is a minimal projective elliptic
surface with exactly two fibres of type $I_0^*$
and $r$ double smooth elliptic fibres $2J_1,\ldots,2J_r$.
Writing $F$ for a full fibre, its actual canonical
line and plurigenera are
\[
 \begin{gathered}
 K_{M_r}\simeq\OO_{M_r}\left(\sum_{j=1}^rJ_j-F\right),
 \qquad
 2K_{M_r}\simeq\OO_{M_r}((r-2)F),
 \\
 P_m(M_r)=\max\{0,\,1-m+r\lfloor m/2\rfloor\}.
 \end{gathered}
                                                               \tag{C.EN.R1}\label{cdp:eq:C.EN.R1}
\]
It has
\[
 (b_0,b_1,b_2,b_3,b_4)=(1,0,10,0,1),\qquad
 (h^0(\OO),h^1(\OO),h^2(\OO))=(1,0,0).
                                                               \tag{C.EN.R2}\label{cdp:eq:C.EN.R2}
\]
The case $r=2$ is the Enriques component of \ref{cdp:st:C.2.UN.25}.
Every $r\ge3$ gives Kodaira dimension one and
coherent Euler characteristic one.

Before balancing, the full central fibre is
\[
 2M_r+2A+4\sum_{j=1}^r Q_j+\sum_{i=1}^4E_i,
                                                               \tag{C.EN.R3}\label{cdp:eq:C.EN.R3}
\]
with $A\simeq\PP^1\times\PP^1$, primary Hopf
$Q_j$, and $E_i\simeq\mathbb F_2$.
Its reduced support has Euler number $12$.
Six of the explicit negative-two operations,
all disjoint from $M_r$, make that number zero.
The reduced supports before and after these
operations have respectively
\[
 \begin{array}{c|c|c}
 &(b_0,b_1,b_2,b_3,b_4)&(h^0(\OO),h^1(\OO),h^2(\OO))\\ \hline
 S_r&(1,1,r+6,2r,r+6)&(1,1,0)\\
 S_{r,\mathrm{bal}}&(1,1,r+18,2r+36,r+18)&(1,1,0).
 \end{array}                                             \tag{C.EN.R4}\label{cdp:eq:C.EN.R4}
\]
All these fillings have the same specified punctured
torus family. Their distinguished essential components
are pairwise non-bimeromorphic, with canonical growth
slope $r/2-1$.
\end{proposition}
\begin{proof}
Use exactly the compact Weierstrass base, quadratic
base change, resolution, and involution of
(C.EN.2)--(C.EN.5). On its exceptional $R_1$,
blow up the $2r$ points
$z=1,\ldots,r,-1,\ldots,-r$.
Replace (C.EN.6) by
\[
 D_r=\sum_{j=1}^r(D_j-D_{-j}),\qquad L_r=\OO(D_r).
                                                               \tag{C.EN.R5}\label{cdp:eq:C.EN.R5}
\]
The canonical section and the involution give
$h^*L_r=L_r^{-1}$ with its actual order-two
section-preserving dualization. The two signed
root charts still have order two, and their
free $q$-quotient and properness are exactly
(C.EN.7)--(C.EN.8). The formula
$\ell=\sigma c^2,\ a=c^{-1}$ supplies an
actual product identification away from the
central divisor, and the involution there
is always $(v,[c])\mapsto(hv,[c^{-1}])$.
It is independent of $r$. This proves the
assertion about the fixed punctured family.

On the strict transform of $R_1$, use
\[
 \phi_r(z)=\prod_{j=1}^r\frac{z-j}{z+j},\qquad
 C_r:\quad Y^2=\prod_{j=1}^r(z^2-j^2),\qquad
 y=\frac{Y}{\prod_{j=1}^r(z+j)}.
                                                               \tag{C.EN.R6}\label{cdp:eq:C.EN.R6}
\]
The $2r$ branch points are distinct, so $C_r$
is smooth and has genus $r-1$, by the
degree-two Riemann--Hurwitz formula.
The actual finite \'etale product cover
of the corresponding component is again
$\ell=u^2,\ a=y/u$ from $C_r\times E$,
with $E=\C^*/q^{\Z}$. Its deck involution
and the lifted dualization are
\[
 \begin{split}
 \eta(z,Y,[u])&=(z,-Y,[-u]),\\
 k(z,Y,[u])&=(-z,(-1)^rY,[u^{-1}]).
 \end{split}                                             \tag{C.EN.R7}\label{cdp:eq:C.EN.R7}
\]
They commute and generate a Klein four group.
The first action is free. The second formula
is forced by the actual identity
$\phi_r(-z)=\phi_r(z)^{-1}$, so it includes
the sign for odd $r$.

We check both parities at the fixed base values.
At infinity, $Y/z^r=\pm1$, and $k_C$
fixes the two points while $jk_C$ exchanges
them. At zero, the two values of $Y$ are
nonzero. When $r$ is even, $k_C$ fixes
them and $jk_C$ exchanges them; when $r$
is odd, the roles are reversed. Therefore
\[
 \begin{array}{c|cc}
 &\#\operatorname{Fix}(k_C)&
       \#\operatorname{Fix}(jk_C)\\ \hline
 r\text{ even}&4&0\\
 r\text{ odd}&2&2 .
 \end{array}                                             \tag{C.EN.R8}\label{cdp:eq:C.EN.R8}
\]
Each nontrivial fixed-point action on the
smooth curve has derivative $-1$.
On $E$, both $[u]\mapsto[u^{-1}]$
and $[u]\mapsto[-u^{-1}]$ have four
fixed points and derivative $-1$.
Consequently the two elements $k,\eta k$
have a total of sixteen fixed points,
all with local representation $(-1,-1)$.
Their sets are disjoint because $\eta$
is free. Dividing by the remaining free
involution pairs them into eight quotient
points. Thus
\[
 X_r=(C_r\times E)/\langle\eta,k\rangle
                                                               \tag{C.EN.R9}\label{cdp:eq:C.EN.R9}
\]
has precisely eight $A_1$ singularities,
and no branch curve. Let $M_r$ be its
minimal $A_1$ resolution. The same
ambient charts (C.EN.14)--(C.EN.15)
realize that exact resolution as a
central component. For odd $r$ the
fixed lifts over zero are the
$\eta k$ lifts; downstairs the local
involution, the four fixed curves,
and all monomials of those charts
are unchanged.

The map $w=z^2$ gives the elliptic map
$p:M_r\to\PP^1_w$. Above a general $w$,
the four points of $C_r$ are acted on
transitively by $\langle\eta,k\rangle$,
so its fibre is exactly $E$. At $w=j^2$,
the local double-root chart and the free
translation in $\eta$ give a smooth
double fibre $2J_j$. At zero and infinity,
the involution has four fixed points on
the elliptic fibre; its $A_1$ resolutions
give the full divisor
$2D+\sum_{i=1}^4 U_i$ of type $I_0^*$,
as in (C.EN.S2). These are all special
values. In particular the morphism $p$
has connected fibres.

We compute its canonical line using an
actual invariant meromorphic form,
instead of deducing it from numerical
invariants. On $C_r\times E$ use
\[
                       \Omega=dz\wedge\frac{du}{u}.
                                                               \tag{C.EN.R10}\label{cdp:eq:C.EN.R10}
\]
It is fixed by both generators in
\eqref{cdp:eq:C.EN.R7}. On $C_r$, $dz$ has a simple
zero at each of the $2r$ branch points
and a double pole at each of the two
points over infinity, and no other
zero or pole. The quotient is free
in codimension one. The root charts
therefore give one zero along every
$J_j$. At infinity, the local invariant
coordinate is $w^{-1}$, whose divisor
pulls back with order two at both
points of $C_r$.

The resolution comparison at the isolated
quotient points retains these orders.
In local coordinates $(s,v)$ where the
involution is $(-s,-v)$, its crepant
volume form satisfies
$ds\wedge dv=\tfrac12\,dr\wedge d\delta$
in the chart $s^2=r,\ sv=r\delta$.
Thus the double pole $s^{-2}$ of
\eqref{cdp:eq:C.EN.R10} becomes a simple pole
on the exceptional curve, and in the
other chart has orders one and two
on the exceptional and strict central
curves. These are exactly the
multiplicities of the full $I_0^*$
fibre. At zero the form is regular
and nowhere zero, and its crepant
extension is likewise nowhere zero.
Consequently the descended form has
the exact divisor
\[
                \operatorname{div}(\Omega)
                         =\sum_{j=1}^r J_j-F_\infty.
                                                               \tag{C.EN.R11}\label{cdp:eq:C.EN.R11}
\]
This proves the first formula of
\eqref{cdp:eq:C.EN.R1}. Since $2J_j$ is linearly
equivalent to a full fibre, it also
proves its second formula.

Every meromorphic section of $mK_{M_r}$
is a meromorphic function times
$\Omega^{\otimes m}$. A holomorphic
such section restricts to a constant
multiple on a general elliptic fibre,
because that fibre's canonical line
is trivial and $\Omega$ has no zero
or pole there. Its ratio is therefore
a function from the connected-fibre
base. The exact allowed orders at
the double fibres in \eqref{cdp:eq:C.EN.R11}
give the sheaf equality
\[
 p_*\OO_{M_r}(mK_{M_r})
  \simeq
 \OO_{\PP^1}\left(
      \sum_{j=1}^r\lfloor m/2\rfloor[j^2]
                       -m[\infty]\right).
                                                               \tag{C.EN.R12}\label{cdp:eq:C.EN.R12}
\]
The ordinary full fibre at infinity,
including all its components, is
Cartier and contributes exactly the
displayed $-m$. Taking sections gives
the plurigenus formula in \eqref{cdp:eq:C.EN.R1}.
For $r=2$, it gives $P_m=1$ for even
$m$ and zero for odd $m$, recovering
\ref{cdp:st:C.2.UN.25}. For $r\ge3$, its linear growth
has positive slope $r/2-1$, so
$\kappa(M_r)=1$. The line $2K_{M_r}$
is a nonnegative multiple of a fibre
and is nef. Adjunction therefore
excludes a $(-1)$-curve, proving that
$M_r$ is minimal for all $r\ge2$.

The component is projective because
it is the resolution of the finite
quotient of the projective surface
$C_r\times E$. Its coherent groups
can also be read through that exact
quotient. On $H^1(C_r,\OO)$ the
hyperelliptic involution $j$ is
$-I$; this follows directly from
the basis $z^a dz/Y$, $0\le a\le r-2$,
of canonical forms and duality.
The action $\eta$ therefore removes
this summand, whereas $k$ removes
the remaining elliptic summand.
Thus $q(M_r)=0$. The $A_1$ resolution
does not change coherent cohomology:
the explicit exceptional $(-2)$-curve
has positive successive conormal
powers, whose $H^1$ vanish, so
formal functions gives $R^1\pi_*\OO=0$.
The formula at $m=1$ gives $p_g=0$.
The finite quotient Euler formula is
$(0+0+16)/4=4$, for both parity
patterns in \eqref{cdp:eq:C.EN.R8}. Resolving
eight $A_1$ points adds eight,
giving $e(M_r)=12$. Hodge symmetry
and Poincar\'e duality on this
projective surface now yield
\eqref{cdp:eq:C.EN.R2}.

The ambient fixed curves over zero
are still the four copies of $R_0$,
and their normal bundles are still
$\OO(-2)\oplus\OO$. The $2r$
root Hopf components are paired
by the involution into $r$ components.
The exact quotient and resolution
monomials give \eqref{cdp:eq:C.EN.R3}, with no
parity change in its coefficients.
The incidence data of \ref{cdp:st:C.2.UN.25} are
unchanged except for the $r$
elliptic attaching curves $M_r\cap Q_j$.
Those curves all have Euler number
zero. Hence the same calculation
gives $e(S_r)=12$, and six \ref{cdp:st:C.2.UN.4}
operations away from $M_r$ give
zero while preserving $M_r$ itself.
The normal crossing and compactness
arguments are the actual charts
of \ref{cdp:st:C.2.UN.25}; no new analytic gluing
is required.

\par\medskip\noindent\textbf{Cohomology of all the paired-root supports.}\par\noindent
The support calculation (C.EN.S1)--(C.EN.S7) extends through
the same maps, with $r$ Hopf components instead of two.
We give the dimensions and the changed maps explicitly.
There are $r+6$ components: the projective elliptic surface
$M$, $A\simeq\PP^1\times\PP^1$, four surfaces
$E_i\simeq\mathbb F_2$, and $Q_1,\ldots,Q_r$.
The double curves are the fourteen rational curves
$D_0,D_\infty,Z_i,U_{i0},U_{i\infty}$ already listed
in (C.EN.S1), and the $r$ elliptic curves
$J_j=M\cap Q_j$. The same eight triple points remain.
On $M$, the actual elliptic fibration gives
\[
 [2D_0+\textstyle\sum_iU_{i0}]
 =[2D_\infty+\textstyle\sum_iU_{i\infty}]
 =2[J_j]\quad(1\le j\le r).
                                                               \tag{C.EN.G1}\label{cdp:eq:C.EN.G1}
\]
Its two $I_0^*$ configurations still span dimension nine
in $H_2(M,\Q)$, by the same affine-$D_4$ intersection
calculation and nonzero fibre class. Here $b_2(M)=10$,
$b_1(M)=b_3(M)=0$, and $h^1(M,\OO_M)=h^2(M,\OO_M)=0$
by the invariant and resolution computations above.

For the exact constant-sheaf complex (C.EN.S3), the
degree-zero row now has dimensions
$r+6,r+14,8$. Its two ranks are $r+5$ and eight:
each extra Hopf vertex and attaching edge is a leaf,
and each of the eight triangles retains its private
edge $U_{i\epsilon}$. Thus its cohomology is still
$(1,1,0)$. In degree one the map has dimensions
$r\to2r$ and rank $r$, since each Hopf radial
generator restricts nontrivially to its own elliptic
attaching curve. In particular this row contributes
dimension $r$ to $H^2$ of the support.

The degree-two restriction map has dimensions
$20\to r+14$ and rank fourteen. The $E_i$ and $A$
equations in (C.EN.S4) are unchanged. The $M$
equation, using \eqref{cdp:eq:C.EN.G1}, becomes
\[
 \sum_{j=1}^r y_j=0,\qquad
 u_{i0}=x_0/2,\quad u_{i\infty}=-x_0/2,\quad
 x_\infty=-x_0,\quad z_i=0.
                                                               \tag{C.EN.G2}\label{cdp:eq:C.EN.G2}
\]
These are all solutions by the already proved
independence of the common half-fibre and the
rational-curve difference subspace. The kernel
therefore has dimension $(r-1)+1=r$, proving
rank fourteen. The component degree-three and
degree-four groups have dimensions $r$ and $r+6$,
respectively. As before, every possible higher
differential has zero target. Consequently
\[
 \bigl(b_i(S_r)\bigr)_{i=0}^4=(1,1,r+6,2r,r+6),
 \qquad e(S_r)=12.
                                                               \tag{C.EN.G3}\label{cdp:eq:C.EN.G3}
\]
For structure sheaves the $r$ Hopf-to-elliptic
restrictions are isomorphisms, and every component
has zero $H^2(\OO)$. The same exact complex gives
$(h^i(S_r,\OO_{S_r}))_{i=0}^2=(1,1,0)$.
Applying six negative blocks away from $M$, with
the support pullback/Gysin maps and derived Cartier
comparison used in (C.EN.S7), gives
\[
 \begin{gathered}
 \bigl(b_i(S_{r,\mathrm{bal}})\bigr)_{i=0}^4
      =(1,1,r+18,2r+36,r+18),\qquad e(S_{r,\mathrm{bal}})=0,\\
 \bigl(h^i(S_{r,\mathrm{bal}},\OO_{S_{r,\mathrm{bal}}})\bigr)_{i=0}^2
      =(1,1,0).
 \end{gathered}                                                \tag{C.EN.G4}\label{cdp:eq:C.EN.G4}
\]
The parameterized incidence and curve-class matrices
are included in \path{code/A_enriques_support.py}; its
executed output checks both the original case and
the first two properly elliptic cases.


Finally a covering family of rational
curves would kill every positive
pluricanonical section by its actual
dominant ruled parametrization.
Thus $M_r$ is non-uniruled, and
\ref{cdp:st:C.2.UN.14} gives its essentiality under
proper birational comparison of
smooth fillings. The canonical
growth slopes in \eqref{cdp:eq:C.EN.R1} are
distinct, so these components
are pairwise non-bimeromorphic.
For $r\ge3$ each is the unique
component of Kodaira dimension
one; the other original components
are rational or Hopf, and the
additional blowup components are
ruled. The same inverse-map
comparison therefore also
distinguishes the corresponding
fillings over the fixed punctured
family.
\end{proof}


